% Slide-deck prelude.
% Mirrors the listing and IVL definitions of ../prelude.tex, minus the
% article-specific packages (enumitem, captions, floats) which clash
% with beamer.  Keep the Move language and the literate table in sync
% with ../prelude.tex.

\usepackage{xcolor}
% Deck palette: one primary blue, used for title bars, boxes, and keywords.
\definecolor{mBlue}{RGB}{28,73,137}
\definecolor{mNavy}{RGB}{16,42,88}   % full-bleed frames: cover and standout
\colorlet{codebg}{mBlue!9}           % frameless tinted background for listings
\usepackage{listings}
\usepackage[T1]{fontenc}
\usepackage[sfdefault,lining]{FiraSans}
\usepackage[scale=0.85,lining]{FiraMono}
% Code in the medium weight of Fira Mono, so listings hold up against Fira Sans.
\newcommand{\codefont}{\ttfamily\fontseries{medium}\selectfont}
\usepackage{relsize}
\usepackage{amsmath,amssymb,array}
\usepackage{esz}
\usepackage{graphicx}
\usepackage{xspace}
\usepackage{tikz}
\usepackage{qrcode}
\usepackage{fontawesome5}
\usetikzlibrary{arrows.meta,positioning,shapes.misc,calc,fit,decorations.pathreplacing}

% Abbreviations
\newcommand{\MVP}[0]{\textsf{MVP}\xspace}
\newcommand{\MSL}[0]{\textsf{MSL}\xspace}
\newcommand{\WP}[0]{\textsf{WP}\xspace}
\newcommand{\solidity}[0]{\textsf{Solidity}\xspace}
\newcommand{\fstar}[0]{\textsf{F*}\xspace}
\newcommand{\verus}[0]{\textsf{Verus}\xspace}
\newcommand{\certora}[0]{\textsf{Certora}\xspace}

% IVL (based on esz.sty)
\newenvironment{ivl}{\small\zedindent=1ex\begin{zed}}{\end{zed}}

\newcommand{\requiresof}[2]{\Zpreop{requires\_of}\langle#1\rangle(#2)}
\newcommand{\abortsof}[2]{\Zpreop{aborts\_of}\langle#1\rangle(#2)}
\newcommand{\ensuresof}[2]{\Zpreop{ensures\_of}\langle#1\rangle(#2)}
\newcommand{\resultof}[2]{\Zpreop{result\_of}\langle#1\rangle(#2)}
\newcommand{\modifiesof}[1]{\Zpreop{modifies\_of}\langle#1\rangle}

\newcommand{\Clo}{\mathcal{C}}
\newcommand{\Par}{\mathcal{P}}
\newcommand{\Fld}{\mathcal{F}}
\newcommand{\Addr}{\Zpreop{address}}
\newcommand{\Mem}{\mathcal{M}}
\newcommand{\Mod}{\Delta}
\newcommand{\Type}{\mathcal{T}}
\newcommand{\Expr}{\mathcal{E}}

\newcommand{\PROC}{\Zkeyword{proc}}
\newcommand{\FUN}{\Zkeyword{fun}}
\newcommand{\AXIOM}{\Zkeyword{axiom}}
\newcommand{\RETURNS}{\Zinrel{returns}}
\newcommand{\HAVOC}{\Zkeyword{havoc}}
\newcommand{\ASSUME}{\Zkeyword{assume}}
\newcommand{\ASSERT}{\Zkeyword{assert}}
\newcommand{\DATATYPE}{\Zkeyword{datatype}}
\newcommand{\IS}{\Zinrel{is}}
\newcommand{\TRUE}{\Zkeyword{true}}
\newcommand{\BOOL}{\Zpreop{bool}}
\newcommand{\CALL}{\Zpreop{call}}

\newcommand{\sat}{\mathrel{{\scriptstyle\vert}\mkern-1mu{\scriptstyle\sim}}}
% Function types in listings are written mathematically: u64 @\ftimes@ u64 @\fto@ u64.
\newcommand{\ftimes}{\ensuremath{\times}}
\newcommand{\fto}{\ensuremath{\rightarrow}}
\newcommand{\slabeled}[2]{#1 \sat #2}
\newcommand{\sq}[1]{\overline{#1}}

% Code
\lstdefinestyle{MoveStyle}{
    basicstyle=\codefont,
    keywordstyle=\color{black},
    commentstyle=\color{gray}\normalfont\itshape,
    escapechar=@,
    breaklines=true,
    literate={|~}{{$\sat$\,}}1
        {<=}{{$\leq$}}2
        {>=}{{$\geq$}}2
        {==}{{$=$}}2
        {!=}{{$\neq$}}2
        {==>}{{$\Longrightarrow$}}3
        {forall}{{$\forall$}}1
        {exists}{{$\exists$}}1,
}

\lstdefinelanguage{Move}{
    morekeywords={
        abort, aborts_if, aborts_of, acquires, address, as, assert, assume,
        borrow_global, borrow_global_mut, break, const, continue, copy,
        copyable, define, drop, else, ensures, ensures_of, exists, false,
        forall, friend, fun, global, has, havoc, if, in, include, invariant,
        key, let, loop, match, modifies, modifies_of, module, move, move_from,
        move_to, mut, native, num, old, onabort, pragma, proof, public,
        reads_of, requires, requires_of, resource, result_of, return, schema,
        script, signer, spec, split, store, struct, true, u8, u64, u128, u256,
        update, use, with, where, while},
    sensitive=true,
    morecomment=[l]{//},
    morecomment=[s]{/*}{*/},
}

% Only | as short inline delimiter; ! would clash with prose on slides.
\lstMakeShortInline[language=Move,style=MoveStyle]|

% Frameless tinted box around display listings: a zero-width frame in the
% background colour so that framesep still pads the text on all sides; the
% x-margins keep the box inside the text width.
\lstdefinestyle{CodeBox}{
    backgroundcolor=\color{codebg},
    frame=single, framerule=0pt, rulecolor=\color{codebg},
    framesep=3pt, xleftmargin=3pt, xrightmargin=3pt,
    aboveskip=\smallskipamount, belowskip=\smallskipamount,
}

% Text between $...$ in a listing is set in bold (for emphasis on slides).
\lstnewenvironment{Move}{
    \lstset{language=Move, style=MoveStyle, style=CodeBox,
        basicstyle=\footnotesize\codefont, keywordstyle=\color{mBlue},
        moredelim=**[is][\bfseries]{\$}{\$}}
}{}

\lstnewenvironment{MoveSmall}{
    \lstset{language=Move, style=MoveStyle, style=CodeBox,
        basicstyle=\scriptsize\codefont, keywordstyle=\color{mBlue},
        moredelim=**[is][\bfseries]{\$}{\$}}
}{}

\lstnewenvironment{MoveDiag}{
    \lstset{style=MoveStyle, style=CodeBox, basicstyle=\scriptsize\codefont}
}{}

% Right-flushed item number at the end of a listing line, e.g. to match an
% invariant to an enumerated bullet: use as @\invnum{1}@ inside a listing.
\newcommand{\invnum}[1]{\hfill{\normalfont\color{mBlue}(#1)}}

% Highlight box for take-aways.
\newcommand{\takeaway}[2][0.9\textwidth]{%
  \par\smallskip
  \centerline{\begin{tikzpicture}
    \node[draw=mBlue, thick, rounded corners, inner sep=5pt,
          text width=#1, fill=mBlue!8]{#2};
  \end{tikzpicture}}\par}
